\chapter{EnvReqs \& Runtime Settings}
\label{ch:envreqs}

\newthought{One block controls the machinery.} \yamlkey{EnvReqs} carries two things: an
entry point for \emph{shared default settings}, so a group of cards can agree on one
runtime profile, and the \yamlkey{V2} map --- the only place scheduling, \file{SAMPLE/}
layout, archiving and the broker are configured. \yamlkey{V2} has exactly ten keys, and
this chapter is all ten.

\section{Only half of this block is checked}
\label{sec:envreqs-zone}

Before the keys, one thing worth knowing, because it decides whether a typo in this block
costs you five seconds or a whole scan: \textbf{a misspelling under \yamlkey{V2} stops the
run; a misspelling anywhere else in \yamlkey{EnvReqs} does not.}

Write \yamlkey{worker\_count} where you meant \yamlkey{workers} and the load fails
immediately, telling you the key you invented and every key you could have meant:

\begin{terminal}
\begin{jout}
unsupported task EnvReqs.V2 setting(s): worker_count;
supported settings are archiver, batch_size, check_modules,
checkpoint_heartbeat_sec, factory, monitor, redis, sample_directory,
worker, workers
\end{jout}
\end{terminal}

Write \yamlkey{Pythonn} instead of \yamlkey{Python} next to it, and nothing happens at
all --- the card loads and the scan runs.

The difference is not an oversight, and the rule behind it is simple: \Jtwo\ checks the
keys it \emph{acts on}, and stays quiet about the keys it merely \emph{records}. Every
setting under \yamlkey{V2} changes how the scan actually executes --- how many Workers,
how results are persisted, which broker --- so a stale or misspelled one must never pass
silently, or you would get a different run than the one you wrote with nothing to tell you
so. The siblings of \yamlkey{V2} --- \yamlkey{OS}, \yamlkey{Python},
\yamlkey{CERN\_ROOT} and any other dependency declaration --- are documentation of what
your project needs. \Jtwo\ preserves them and acts on none of them, so it has no standing
to judge how you spelled them.

The practical consequence is worth carrying into the rest of the chapter: when a
\yamlkey{V2} setting does not seem to be taking effect, the problem is not a typo. A typo
would have stopped you at load time. Look instead at the shared defaults file below, or at
the short list of keys that are accepted and then deliberately ignored
(Sections~\ref{sec:sample-directory} and~\ref{sec:archiver-policy}).

\section{Shared defaults: \yamlkey{Check\_default\_dependencies}}
\label{sec:check-default-deps}

Projects that keep a common environment file can load runtime defaults from it:

\begin{yamlwin}
EnvReqs:
  Check_default_dependencies:
    required: true
    default_yaml_path: "&J/deps/environment_default.yaml"
\end{yamlwin}

When \yamlkey{required} is true, the loader resolves \yamlkey{default\_yaml\_path} (normal
\marker{\&J/} rules), loads that file, and reads \emph{only} its \yamlkey{EnvReqs.V2} map.
Those defaults are recursively merged \emph{under} the task's own \yamlkey{EnvReqs.V2} ---
the task always wins, key by key, so a card can override one setting without restating the
profile. With \yamlkey{required: false} the file is skipped; a required but missing or
malformed file fails the load with a focused error.

The shared file is held to the same standard as the card: an unsupported key in
\emph{its} \yamlkey{EnvReqs.V2} is the same hard error, reported against the defaults
file rather than the card.

\section{The \yamlkey{V2} map at a glance}
\label{sec:v2-map}

\begin{yamlwin}
EnvReqs:
  V2:
    workers: 4              # Worker processes
    batch_size: 256         # sampler submit group
    sample_directory:       # SAMPLE/ bucket policy
      limit: 200
    archiver:               # persistence policy
      mode: process
      batch_size: 200
    worker:                 # per-Worker policy
      force_serial_layers: false
      sample_artifacts: auto
    factory:                # supervision
      watchdog: {enabled: true, stale_sec: 30.0}
    monitor:                # monitor snapshot rate
      hz: 120
    redis:                  # broker location
      host: 127.0.0.1
      port: 6379
    check_modules:          # `Jarvis check` budget
      n_samples: 10
    checkpoint_heartbeat_sec: 30   # checkpoint cadence
\end{yamlwin}

Every one of those keys is optional; the card above is a complete enumeration, not a
recommendation. Most production cards set \yamlkey{workers} and nothing else.

\section{Scheduling: \yamlkey{workers} and \yamlkey{batch\_size}}
\label{sec:workers}

\begin{keylist}
  \item[workers] \opt, integer $\geq 0$, default \code{0}. The number of long-lived
        Worker processes. Zero or negative means \emph{one} Worker.
  \item[batch\_size] \opt, integer $> 0$, default \code{256}. How many task
        dictionaries the sampler pushes to Redis per submit group. This is queue
        pipelining only --- execution is always one Sample per Worker --- so the default
        is almost always right.
\end{keylist}

Sizing \yamlkey{workers} is the single most important performance decision you make;
Chapter~\ref{ch:performance} gives the full guidance. The short version: for
calculator-bound scans, set it near the number of calculator invocations your machine (and
license pool) can genuinely sustain.

\paragraph{\yamlkey{worker}, singular, is not the count.} The two spellings are
different settings, and the missing \code{s} is the single easiest slip to make in this
block. \yamlkey{workers} is the count; \yamlkey{worker} is the per-Worker execution
policy \emph{map} of Section~\ref{sec:worker-factory}. Writing a number where the map
belongs stops the load rather than quietly running one Worker:

\begin{terminal}
\begin{jout}
EnvReqs.V2.worker is a policy mapping; use workers: N for the worker count
\end{jout}
\end{terminal}

The two may sit side by side --- \yamlkey{workers: N} alongside
\yamlkey{worker: \{\ldots\}} is the normal way to set a count and a policy at once.

\begin{table}[ht]
  \footnotesize
  \begingroup
  \setlength{\tabcolsep}{0pt}
  \begin{tabular}{@{}p{0.255\linewidth}p{0.213\linewidth}p{0.532\linewidth}@{}}
    \toprule
    \textbf{Key} & \textbf{Default} & \textbf{Meaning} \\
    \midrule
    \yamlkey{enabled} & \code{true} & allocate Samples to buckets on pull \\
    \yamlkey{limit}   & \code{200}  & Samples per bucket directory \\
    \yamlkey{width}   & \code{6}    & zero-pad width of bucket names (\file{000001}) \\
    \yamlkey{pack}    & \code{true} & tar a bucket once fully archived \\
    \yamlkey{start\_bucket} & \code{1} & first bucket number \\
    \bottomrule
  \end{tabular}
  \endgroup
  \caption{\yamlkey{sample\_directory} keys.}
  \label{tab:sample-directory}
\end{table}

\section{\yamlkey{sample\_directory}: the SAMPLE bucket policy}
\label{sec:sample-directory}

Controls the bucketed layout introduced in Section~\ref{sec:sample-lifecycle}:
\file{SAMPLE/\allowbreak{}<bucket>/\allowbreak{}<uuid>/}, packed to
\file{<bucket>.\allowbreak{}tar.\allowbreak{}gz} after archiving. Its keys are in
Table~\ref{tab:sample-directory}.

A bucket is packed only when it is sealed, no Sample in it is still active, and every
Sample in it has been archived to the database --- packing can never outrun persistence.

\paragraph{Where to write it.} \yamlkey{EnvReqs.V2} is the only home for this map ---
\yamlkey{Scan} has no key of its own for it (\S\ref{sec:scan-block}), so there is nowhere
else it could be shadowed from.

\begin{jwarn}
Two spellings, \yamlkey{archive\_samples} and \yamlkey{archive\_prune}, are accepted here
and then ignored --- they pass validation and change nothing. Bucket packing is controlled
by \yamlkey{pack} above and by the Archiver's \yamlkey{pack\_buckets}
(Section~\ref{sec:archiver-policy}).
\end{jwarn}

\begin{table}[ht]
  \footnotesize
  \begingroup
  \setlength{\tabcolsep}{0pt}
  \begin{tabular}{@{}p{0.255\linewidth}p{0.213\linewidth}p{0.532\linewidth}@{}}
    \toprule
    \textbf{Key} & \textbf{Default} & \textbf{Meaning} \\
    \midrule
    \yamlkey{mode} & \code{process} & dedicated Archiver process, or a thread in the
                                       control process (\code{thread}) \\
    \yamlkey{pack\_buckets} & \code{true} & tar sealed buckets after archiving \\
    \yamlkey{batch\_size} & \code{200} & Samples buffered per database flush \\
    \yamlkey{flush\_interval\_sec} & \code{1.0} & time-based flush for partial batches
                                                  (floor \code{0.05}) \\
    \yamlkey{max\_hdf5\_bytes} & \code{1073741824} & shard the database once the current
                                       \file{*.hdf5} passes this size (1\,GiB) \\
    \yamlkey{strategy} & \code{move} & accepted, inert (see below) \\
    \yamlkey{delete\_after\_archive} & \code{true} & accepted, inert (see below) \\
    \bottomrule
  \end{tabular}
  \endgroup
  \caption{\yamlkey{archiver} keys.}
  \label{tab:archiver-keys}
\end{table}

\section{\yamlkey{archiver}: persistence policy}
\label{sec:archiver-policy}

Table~\ref{tab:archiver-keys} lists the keys that decide how much of a finished Sample
survives on disk.

The defaults --- a dedicated process, batched flushes, packed buckets --- are the
production configuration; Chapter~\ref{ch:distributed-runtime} explains the machinery
behind them and when to deviate.

\begin{jwarn}
\yamlkey{strategy} and \yamlkey{delete\_after\_archive} are accepted and then ignored.
A Worker writes each Sample's products straight to their final home under \file{SAMPLE/}
(Section~\ref{sec:sample-lifecycle}), so there is no intermediate copy for either key to
govern --- setting them changes nothing. The same is true of \yamlkey{archive\_samples}
and \yamlkey{archive\_prune} under \yamlkey{sample\_directory}
(Section~\ref{sec:sample-directory}).
\end{jwarn}

\section{\yamlkey{worker}, \yamlkey{factory}, \yamlkey{monitor}: execution and supervision}
\label{sec:worker-factory}

Three small policy groups you will rarely set, and should recognise when you meet them in
somebody else's card. Each takes a map, and each accepts exactly the keys below --- an
invented one is rejected by name, the same as at the top level of \yamlkey{V2}.

\begin{keylist}
  \item[worker.force\_serial\_layers] \opt, default \code{false}. Run a Sample's
        execution plan strictly one calculator at a time, even where the declared
        dependencies would allow concurrency (Chapter~\ref{ch:workflow-model}). A
        debugging tool: it makes an interleaved failure reproducible, at the cost of the
        second concurrency axis.
  \item[worker.sample\_artifacts] \opt, one of \code{auto} (default), \code{always},
        \code{never}. Whether a Sample's working directory is materialised on disk.
        \code{auto} creates it only when something actually needs it --- a
        \token{@Sdir} reference, a file-based calculator --- so opera-only scans write
        nothing per point.
  \item[factory.watchdog] \opt, and \yamlkey{watchdog} is the \emph{only} key
        \yamlkey{factory} takes. A map: \yamlkey{enabled} (\code{true}),
        \yamlkey{stale\_sec} (\code{30.0}), \yamlkey{poll\_interval\_sec} (\code{1.0}),
        \yamlkey{max\_sample\_retries} (\code{3}). Detects a Worker that has stopped
        reporting and re-queues the Sample it was holding, up to the retry limit.
  \item[monitor.hz] \opt, default \code{120}, floor \code{1}, and likewise the only key
        \yamlkey{monitor} takes. How often the \code{TaskFactory} refreshes the snapshot
        \cli{Jarvis monitor} reads (Chapter~\ref{ch:monitoring}). Note the home: this is
        a sibling of \yamlkey{factory}, not a key inside it.
\end{keylist}

\section{\yamlkey{redis}: where the broker lives}
\label{sec:redis-settings}

\begin{keylist}
  \item[host] \opt, default \code{127.0.0.1}.
  \item[port] \opt, default \code{6379}; must be in \code{1}--\code{65535}, and an
        out-of-range or unparseable value falls back to the default rather than failing.
  \item[db] \opt, default \code{0}.
\end{keylist}

Those three are the only keys read; anything else under \yamlkey{redis} is ignored. The
defaults are the zero-configuration local broker \Jtwo\ starts and manages for you, which
is what you want on a laptop or a single node. Set them when the broker is somewhere
else --- a shared instance, a non-default port because something else already owns 6379,
or a separate database number to keep two projects' keys apart.

\begin{jnote}
Changing \yamlkey{port} changes which broker the scan uses, and therefore which scan
\cli{Jarvis monitor} and \cli{Jarvis ps} attach to --- they discover a running scan
through the runtime metadata it published, so they follow the card rather than assuming
6379 (Section~\ref{sec:cli-ps-kill}).
\end{jnote}

\section{\yamlkey{check\_modules}: the smoke-test budget}
\label{sec:check-modules-settings}

Three keys, consulted only by \cli{Jarvis check} (Chapter~\ref{ch:check-modules}):

\begin{keylist}
  \item[data] \opt, default \marker{\&J/}\file{data/\allowbreak{}check\_modules\_points.csv}.
        A \textsc{csv} of fixed points to rehearse with, if that file exists. This is the
        \emph{only} place a check \textsc{csv} can be named.
  \item[n\_samples] \opt, default \code{10}. How many points to draw from your real
        \yamlkey{Sampling.Method} when no \textsc{csv} resolves.
  \item[timeout\_sec] \opt, default \code{120.0}. How long \cli{check} waits for its
        smoke Samples to reach the database --- a deadline, not a fixed sleep, so a fast
        card finishes fast. \cli{Jarvis check --timeout} overrides it for one invocation
        (\S\ref{sec:cli-check}).
\end{keylist}

\section{\yamlkey{checkpoint\_heartbeat\_sec}: how often state is saved}
\label{sec:checkpoint-heartbeat}

\opt, a positive number, default \code{30}, floor \code{1}. The cadence at which a
sampler that snapshots \emph{during} a run writes its state: the fixed-set and
\sampler{CSV} methods (Chapters~\ref{ch:random}--\ref{ch:csv-sampler}) and the
\textsc{mcmc} family (Chapters~\ref{ch:toymcmc}--\ref{ch:ptmcmc}).

It sets a floor, not a timer. The snapshot is taken at the first point the sampler can
safely take one once this many seconds have elapsed --- a feedback boundary, never
mid-generation --- so a checkpoint always describes a consistent state
(Chapter~\ref{ch:checkpoint}). Lower it when a run is expensive to lose and cheap to
snapshot; raise it when the state is large.

Nested sampling does not use this value: \sampler{MultiNest} and \sampler{Dynesty} keep
their own engine-save cadence (Section~\ref{sec:nested-distributed}), so changing it has
no effect on those two.

\section{What is \emph{not} here}
\label{sec:not-here}

Deliberate omissions, so you do not go looking for them:

\begin{itemize}
  \item \textbf{Per-sampler thread counts or Worker pinning.} Workers compete on one
        shared queue; there is nothing to pin.
  \item \textbf{A mode switch.} There is exactly one runtime path
        (Section~\ref{sec:runtime-picture}) --- a twenty-point test and a million-point
        production scan take the same code path.
  \item \textbf{Per-calculator concurrency limits.} Those are pools, and they belong to
        the calculator that needs them (Chapter~\ref{ch:calculators}).
  \item \textbf{Output paths.} The output root follows the project and the scan name
        (Section~\ref{sec:scan-block}); nothing under \yamlkey{V2} moves it.
\end{itemize}
